\documentclass{article}[12pt]

\usepackage{times}
\usepackage{indentfirst}
\usepackage{hyperref}

\author{Violet French}
\title{Fullstack Development Final Technical Document}
\date{\today}

\begin{document}
    \maketitle
    \newpage
    \tableofcontents

    \section{Executive Summary}
    \subsection{What is a Fullstack App}
    A Fullstack appliaction is an application that brings together the fullstack of technologies those being a frontend, backend, and database. The frontend is the part of the app that the client user will see and is typically written in HTML, CSS, and JS. In some cases these files are dynamically rendered on the server side and delivered to the client on request. In our case, we did this with a framework called React.
    
    \subsection{Technical Progression}
    Throughout the semester we worked with a broad range of technologies. We started the semester working with HTML and CSS building static webpages. After we got the hang of the structure of webpages we moved onto JavaScript to make our webpages interactable. After becoming familiar with JavaScript we moved onto a JavaScript framework called React to make building our webpages easier. We used React to build the frontend of our final application. Once we finished learning the frontend we moved onto the backend. For the backend we worked with Node JS which is a JavaScript runtime environment. Using Node's package manager we installed a framework called Express to make building the backend routes easier. The routes would respond to requests from our frontend. Finally we set up a no-SQL database called MongoDB and used the Mongoose package on our backend to communicate with the database. 
    
    \subsection{Skill Progression}
    Throughout the semester I improved in many facets. I got much better at working with CSS, a skill I have remained deficient in for a long time. My JavaScript abilities also improved as I got more familiar with filtering and input validation. I also improved a lot in my abilities to document code. With Express I had to build my own structure for documentation since Express documentation remains largely unspecified. I also improved in my abilities to write \LaTeX which is what I used for much of my external documentation. I learned how to better work with React, a framework I had only breifly worked with in the past. Lastly I learned MongoDB and what no-SQL databases look like. In the past I have always worked with SQL based databases and it was a nice change of pace to work with Mongo which was very easy to implement into JavaScript.

    \section{Frontend Development Foundations}
    \subsection{Using HTML, CSS, and JS}
    At the start of the semeseter we worked with just HTML, CSS, and JS. We started by just working with HTML and CSS to build basic webpages that users couldn't interact with in any way. With HTML we discussed how to create elements using HTML tags. We also discussed organizing our document using semantic elements like sections or articles. Using semantic elements and logical structures helps make our site more accessible and assists screenreaders in parsing our webpages. With CSS we talked about how to style different elements. We discussed flexboxes and gridboxes as ways to organize our HTML elements on screen in visually appealing ways. We discussed media tags which are a way to make our webpages responsive. Responsive webpages will restyle themselves depending on the system or display they are being rendered on. They may make elements bigger on smaller screens to make them easier to see on phones or tablets. This improves upon the older method of hosting multiple sites for different devices. Finally with JS we discussed how to make our webpages interactable. We worked with JS arrays to store collections of information. We also worked with various array functions allowing us to filter or map elements in arrays. We worked with JavaScript objects to store various related properties in one location. We used functions to section off our code into reusable fragments. These functions could then be executed when an event, such as pressing a button or changing a form input, occurs. Finally we discussed DOM manipulation so we could create new HTML elements or reference certain attributes of existing elements. This allowed us to pull information from inputs or create new elements to display to the user.
    
    \subsection{React}
    After discussing the basics of HTML, CSS, and JS we moved onto the more streamlined React framework. React took the cumbersome DOM manipulation of JavaScript and made it far easier to work with. Instead of having to manually rerender elements anytime some value changed React handles those changes for us. React does this through state variables. When a state variable's setter function is called it will change the state's value and then rerender any components that rely on it. React also provides us with JSX an HTML like markup langauge that can be used to create lightweight component for rendering. Rather than having our HTML and JS in sepereate files React stores it all in one file using components. Each React component will be passed properties like states and will return JSX describing how to build that component. The component may also have function inside it for handling things like button presses or form submissions.
    
    \subsection{Vite}
    React requires a special compiler to convert it to HTML, CSS, and JS that the clients browser can read. Additionally a server is needed to send these webpages to the user upon request. Vite handles all this for us. Vite is a frontend build tool and development server for JavaScript and TypeScript. Vite will bundle together all our modules, compile our React code, and serve it up to us at a specific host and port. Vite allows us to view our webpage as we build it using Hot Module Replacement which will rerender parts of the webpage as they change. This accelerates development and allows us to test our page before it is deployed to the wider internet.

    \subsection{Technical Evidence}

    %%%%%%%% TODO %%%%%%%%%

    \section{React Architecture and State Management}
    \subsection{Overview of React}
    \subsubsection{Components}
    
    \subsubsection{Props}
    
    \subsubsection{State}
    % Include part about updating arrays and objects without mutating state
    % https://react.dev/learn/updating-arrays-in-state

    \subsubsection{Effects}
    
    \subsubsection{Control Forms}
    
    \subsubsection{Conditional Rendering}
    
    % Show how data moves with example
    % Include a small bit of code demonstrating state or component communication and explain it

    \section{Backend Develeopment and REST}
    \subsection{Node}

    \subsection{Express and HTTP}
    
    \subsubsection{Routes and Middleware}
    
    \subsubsection{HTTP Methods and Response Code}
    
    \subsubsection{Request Params and Bodies}
    
    \subsubsection{Cross Origin Resource Sharing}
    
    \subsubsection{Thunderclient}
    
    \subsubsection{Error Handling and Validation}

    % Trace one GearVault operation such as creating or editing equipment from the React interface through express to the database and back

    \section{Database Integration and Persistence}
    \subsection{MongoDB}
    
    \subsection{Mongoose}
    
    \subsection{Mongo Concepts}
    
    \subsubsection{Collections and Documents}

    % Discuss _id
    
    \subsubsection{Schemas and Models}
    
    \subsubsection{Transactions}
    
    \subsection{Local Storage vs Database Persistence}
    
    \subsection{Equipment and Reservations in GearVault}

    % Explain how GearVault prevents equipment item from being reserved by two people at the same time including the role of backend availability checks
    % The availability checks ensure that we only try to update things with valid equipment ids that are available

    \section{Fullstack Architecture and Integration}

    % Insert Architectural Diagram showing how the whole GearVault app works
    % Should include example of the architecture layers with explanations
        % Diagram should be annotated
    % Explain how data moves through the system when a user reserves equipment and when it's returned
    % Identify the frontend, backend, and databases' responsabilities

    \section{Testing, Debugging, and Problem Solving}
    \subsection{Problem 1}
    \subsubsection{Symptoms}
    
    \subsubsection{Investigation}

    \subsubsection{Cause}

    \subsubsection{Solution}

    \subsubsection{Reflection}

    \subsection{Problem 2}
    \subsubsection{Symptoms}
    
    \subsubsection{Investigation}

    \subsubsection{Cause}

    \subsubsection{Solution}

    \subsubsection{Reflection}
    
    \subsection{Problem 3}    
    \subsubsection{Symptoms}
    
    \subsubsection{Investigation}

    \subsubsection{Cause}

    \subsubsection{Solution}

    \subsubsection{Reflection}

    \section{Personal Reflection and Future Improvements}
    \subsection{Biggest Challenge}
    
    \subsection{Most Important Skill Developed}

    \subsection{Improvements to Gear Vault}
    % Explain the technical value of each proposed improvement
    \subsubsection{Improvement 1}
    
    \subsubsection{Improvement 2}


\end{document}